\section*{Каноническая проекция ротора на фиксированный импульс}

Для свободного квантового ротора
\[
 H=\frac{p^2}{2I},\qquad p\in\mathbb Z,
\]
вклад фиксированного импульсного сектора равен
\[
 Z_p=\exp\!\left(-\frac{\beta p^2}{2I}\right).
\]
Он содержит обратную инерцию, но не содержит независимого члена
\(\log I\). Корневой множитель гауссовых флуктуаций появляется в сумме по
намоткам и сокращается при пуассоновском переходе к отдельным импульсным
секторам. Поэтому один канонический ротор не может одновременно породить
логарифмический и обратный члены функционала щели.